\documentclass[12pt,a4paper]{article}

\usepackage[T2A]{fontenc}
\usepackage[utf8]{inputenc}
\usepackage[russian]{babel}

\usepackage{amsmath,amssymb,amsfonts}
\usepackage{mathtools}
\usepackage{graphicx}
\usepackage{float}
\usepackage{booktabs}
\usepackage{longtable}
\usepackage{array}
\usepackage{caption}
\usepackage{subcaption}
\usepackage{multirow}
\usepackage{tabularx}
\usepackage{makecell}
\usepackage{tocloft}
\usepackage{xcolor}
\usepackage{geometry}
\usepackage{hyperref}

\geometry{
    left=25mm,
    right=25mm,
    top=25mm,
    bottom=25mm
}

\hypersetup{
    colorlinks=true,
    linkcolor=blue!50!black,
    urlcolor=blue!50!black,
    citecolor=blue!50!black,
    pdftitle={СР_1 по анализу временных рядов},
    pdfauthor={Ковалёва Мария, БЭК2310}
}

\setlength{\parindent}{0pt}
\setlength{\parskip}{0.4em}
\setlength{\tabcolsep}{4pt}
\emergencystretch=1.5em
\renewcommand{\arraystretch}{1.12}
\renewcommand{\cftsecleader}{\cftdotfill{\cftdotsep}}
\renewcommand{\cftsubsecleader}{\cftdotfill{\cftdotsep}}

\newcommand{\placeholder}{\textit{Раздел будет заполнен на следующем этапе.}}
\renewcommand\cellalign{tl}

\title{Самостоятельная работа 1\\Анализ временных рядов занятости в сфере услуг для Германии и Турции}
\author{Ковалёва Мария, БЭК2310}
\date{\today}

\begin{document}

\maketitle
\thispagestyle{empty}
\clearpage
\tableofcontents
\newpage

\section{Введение}

Доля занятых в сфере услуг является содержательно важным индикатором структурной трансформации экономики, поскольку отражает перераспределение занятости между секторами и степень перехода к более сервисно-ориентированной модели хозяйства. Рост этого показателя обычно связан с изменением отраслевой структуры, урбанизацией, развитием торговли, финансовых и иных услуг, а также с долгосрочными сдвигами в производительности и потребительском спросе. Поэтому анализ динамики показателя \textit{Employment in services (\% of total employment)} позволяет исследовать не только изменение занятости как таковой, но и более широкий процесс экономической модернизации.

В работе рассматриваются два временных ряда для Германии и Турции по данным World Bank \url{https://data.worldbank.org/indicator/SL.SRV.EMPL.ZS}. Выбор именно этих стран обусловлен тем, что они представляют разные типы экономической динамики. Германия выступает как зрелая развитая экономика с высоким уровнем сервисизации и более плавной траекторией изменений, тогда как Турция характеризуется более выраженной догоняющей динамикой и более интенсивным переходом к сервисной структуре занятости. Такое сопоставление позволяет содержательно сравнить не только уровни показателя, но и различия в скорости и характере структурных сдвигов.

Период наблюдения 1991--2025 гг. выбран по двум причинам. Во-первых, он обеспечивает достаточно длинный годовой ряд для применения методов анализа временных рядов и построения ARIMA-моделей. Во-вторых, данный интервал охватывает как длительные структурные изменения, так и кризисные эпизоды и современный этап развития экономик, что делает возможным исследование тренда, потенциальной нестационарности и структурных сдвигов в сопоставимой для обеих стран временной рамке.

Актуальность исследования определяется тем, что анализ структурных изменений занятости и динамики сервисного сектора является важной задачей для понимания долгосрочного экономического развития стран. Использование методов анализа временных рядов позволяет формализовать эти процессы и получить количественные оценки трендов и прогнозов.

В рамках работы последовательно выполняются описание исходных данных и анализ динамики показателя, исследование трендовой компоненты и структурных сдвигов, определение порядка интегрируемости процесса, оценивание и сравнение нескольких спецификаций ARIMA-моделей, анализ остатков выбранных моделей, а также построение прогноза и оценка его качества. Такой порядок анализа позволяет перейти от содержательного описания рядов к формализованному моделированию и сопоставить особенности динамики Германии и Турции на единой методологической основе.

\section{Исходные данные}

В данном разделе приводится краткое описание исходных данных, используемых для дальнейшего сравнительного анализа временных рядов. Основные характеристики выборки сведены в таблицу~\ref{tab:data_description}.

\begin{table}[H]
\centering
\begin{flushright}
\small Таблица 1
\end{flushright}
\caption{Характеристика исходных данных по показателю доли занятых в сфере услуг, \%, 1991--2025 гг.}
\label{tab:data_description}
\begin{tabularx}{\textwidth}{>{\raggedright\arraybackslash}p{0.34\textwidth} >{\raggedright\arraybackslash}X}
\toprule
\textbf{Характеристика} & \textbf{Значение} \\
\midrule
Анализируемый показатель & Employment in services (\% of total employment) \\
Единица измерения & \% \\
Страны & Германия, Турция \\
Источник данных & World Bank, World Development Indicators, \url{https://data.worldbank.org/indicator/SL.SRV.EMPL.ZS} \\
Период наблюдения & 1991--2025 \\
Количество наблюдений & 35 на каждую страну, 70 всего \\
\bottomrule
\end{tabularx}
\end{table}

Динамика показателя по странам представлена на рис.~\ref{fig:combined_series_section2}. На общем графике видно, что уровень показателя на всём интервале выше в Германии, тогда как траектория Турции имеет более крутой наклон и тем самым отражает более быстрый рост доли занятых в сфере услуг.

\begin{figure}[H]
\centering
\includegraphics[width=0.92\textwidth]{figures/combined_series_section2.png}
\caption{Сравнение динамики показателя Employment in services (\% of total employment) для Германии и Турции, \% от общей занятости, 1991--2025 гг.}
\label{fig:combined_series_section2}
\end{figure}

За рассматриваемый период в Германии показатель увеличился с 57.92\% в 1991 г. до 72.95\% в 2025 г., то есть на 15.03 процентного пункта, или на 25.95\% относительно начального уровня. В Турции рост оказался существенно более быстрым: показатель повысился с 31.13\% до 58.49\%, что соответствует увеличению на 27.36 процентного пункта, или на 87.89\%. Таким образом, Германия сохраняет более высокий уровень сервисной занятости, а Турция демонстрирует более интенсивное повышение показателя.

\begin{table}[H]
\centering
\begin{flushright}
\small Таблица 2
\end{flushright}
\caption{Дескриптивные статистики показателя доли занятых в сфере услуг, \%, 1991--2025 гг.}
\label{tab:desc_stats}
\begin{tabular}{lrrrrrr}
\toprule
\textbf{Страна} & \textbf{count} & \textbf{mean} & \textbf{median} & \textbf{std} & \textbf{min} & \textbf{max} \\
\midrule
Германия & 35 & 67.31 & 69.04 & 4.61 & 57.92 & 72.95 \\
Турция & 35 & 46.03 & 48.29 & 8.45 & 31.13 & 58.49 \\
\bottomrule
\end{tabular}
\end{table}

Таблица~\ref{tab:desc_stats} показывает, что средний уровень показателя существенно выше в Германии: 67.31\% против 46.03\% в Турции. Разброс значений, напротив, больше у Турции, где стандартное отклонение достигает 8.45, тогда как в Германии оно составляет 4.61. По минимальным и максимальным значениям Германия также находится выше на всём интервале наблюдения: диапазон составляет от 57.92\% до 72.95\%, тогда как для Турции он лежит в пределах от 31.13\% до 58.49\%. Следовательно, Германия характеризуется более высоким и более ровным уровнем занятости в сфере услуг, а Турция --- более низкой стартовой базой и большей вариативностью динамики.

В целом сопоставление двух рядов выявляет устойчивое различие как по уровню, так и по характеру изменения показателя. Германия отражает более зрелую структуру занятости с плавной динамикой, тогда как Турция демонстрирует более интенсивное повышение доли занятых в сфере услуг. Именно сочетание различий в уровне, темпе роста и вариативности делает эти ряды содержательно интересными для дальнейшего анализа тренда, стационарности, структурных сдвигов и ARIMA-моделирования.

\section{Анализ временного ряда}

\subsection{Основные компоненты временного ряда}

Сопоставление рядов показывает, что в обоих случаях наблюдается устойчивый восходящий тренд, однако его форма различается. Для Германии характерен более плавный и ровный рост доли занятых в сфере услуг, тогда как для Турции рост выражен сильнее и идет заметно быстрее. Это различие видно как по исходной динамике, так и по сглаженным траекториям на рис.~\ref{fig:germany_trend_section3} и~\ref{fig:turkiye_trend_section3}: германский ряд выглядит более стабильным, а турецкий --- более интенсивным по наклону и локальным изменениям.

\begin{figure}[H]
\centering
\includegraphics[width=0.88\textwidth]{figures/germany_trend_section3.png}
\caption{Трендовая динамика показателя Employment in services (\% of total employment) для Германии}
\label{fig:germany_trend_section3}
\end{figure}

\begin{figure}[H]
\centering
\includegraphics[width=0.88\textwidth]{figures/turkiye_trend_section3.png}
\caption{Трендовая динамика показателя Employment in services (\% of total employment) для Турции}
\label{fig:turkiye_trend_section3}
\end{figure}

Коррелограммы исходных рядов приведены на рис.~\ref{fig:acf_germany_section3}, \ref{fig:pacf_germany_section3}, \ref{fig:acf_turkiye_section3} и~\ref{fig:pacf_turkiye_section3}. Для обеих стран автокорреляционная функция убывает медленно, а не обрывается после первых лагов, что указывает на сильную инерционность рядов и позволяет предположить нестационарность уровня. Такая форма ACF согласуется с наличием выраженной трендовой компоненты: текущие значения существенно зависят от предшествующей траектории, а не колеблются вокруг постоянного среднего. При этом для Турции затухание воспринимается как более устойчивое, что соответствует более интенсивному росту показателя по сравнению с Германией.

\begin{figure}[H]
\centering
\includegraphics[width=0.70\textwidth]{figures/acf_germany_section3.png}
\caption{ACF исходного ряда для Германии}
\label{fig:acf_germany_section3}
\end{figure}

\begin{figure}[H]
\centering
\includegraphics[width=0.70\textwidth]{figures/pacf_germany_section3.png}
\caption{PACF исходного ряда для Германии}
\label{fig:pacf_germany_section3}
\end{figure}

\begin{figure}[H]
\centering
\includegraphics[width=0.70\textwidth]{figures/acf_turkiye_section3.png}
\caption{ACF исходного ряда для Турции}
\label{fig:acf_turkiye_section3}
\end{figure}

\begin{figure}[H]
\centering
\includegraphics[width=0.70\textwidth]{figures/pacf_turkiye_section3.png}
\caption{PACF исходного ряда для Турции}
\label{fig:pacf_turkiye_section3}
\end{figure}

Частные автокорреляционные функции дополняют этот вывод. На рис.~\ref{fig:pacf_germany_section3} и~\ref{fig:pacf_turkiye_section3} в обоих рядах выделяется прежде всего первый лаг, после которого частные автокорреляции заметно ослабевают и становятся менее устойчивыми. Это означает, что существенная зависимость концентрируется на малом числе первых лагов, а дальнейшая корреляционная структура во многом связана с общей трендовой динамикой. В сравнительном плане PACF Турции выглядит несколько более выраженной по первым лагам, что согласуется с более быстрым переходом к сервисной занятости. Следовательно, уже на этом этапе можно сделать два вывода: оба ряда содержат трендовую компоненту и, вероятно, нестационарны в уровнях, а возможная AR-структура, если переходить к моделированию, должна проверяться прежде всего в низких порядках.

\subsection{Проверка гипотезы об отсутствии тренда}

Для формальной проверки гипотезы об отсутствии тренда использованы два теста, не связанные с проверкой единичного корня: критерий серий относительно медианы (\textit{Runs test}) и тест \textit{Kendall tau} на монотонный тренд. В первом случае нулевая гипотеза состоит в том, что последовательность наблюдений случайна относительно медианы и систематического тренда нет. Во втором случае нулевая гипотеза формулируется как отсутствие монотонной связи между временем и уровнем ряда. Результаты тестирования приведены в таблице~\ref{tab:trend_tests_section3}.

\begin{table}[H]
\centering
\begin{flushright}
\small Таблица 3
\end{flushright}
\caption{Проверка гипотезы об отсутствии тренда для показателя доли занятых в сфере услуг, \%, 1991--2025 гг.}
\label{tab:trend_tests_section3}
\scriptsize
\begin{tabularx}{\textwidth}{@{}>{\raggedright\arraybackslash}p{0.11\textwidth} >{\raggedright\arraybackslash}p{0.08\textwidth} >{\raggedright\arraybackslash}p{0.24\textwidth} >{\raggedleft\arraybackslash}p{0.10\textwidth} >{\raggedleft\arraybackslash}p{0.08\textwidth} >{\raggedright\arraybackslash}X@{}}
\toprule
\makecell[tl]{\textbf{Страна}} & \makecell[tl]{\textbf{Тест}} & \makecell[tl]{\textbf{H0}} & \makecell[tr]{\textbf{Статис-} \\ \textbf{тика}} & \makecell[tr]{\textbf{p-} \\ \textbf{value}} & \makecell[tl]{\textbf{Вывод}} \\
\midrule
Германия & Runs & Последовательность случайна относительно медианы, систематического тренда нет. & -5.6618 & $< 0.001$ & H0 отвергается: порядок наблюдений неслучаен, тренд вероятен. \\
Германия & Kendall & Монотонной связи между временем и уровнем ряда нет. & 0.9664 & $< 0.001$ & H0 отвергается: присутствует статистически значимый монотонный тренд. \\
Турция & Runs & Последовательность случайна относительно медианы, систематического тренда нет. & -4.9749 & $< 0.001$ & H0 отвергается: порядок наблюдений неслучаен, тренд вероятен. \\
Турция & Kendall & Монотонной связи между временем и уровнем ряда нет. & 0.9462 & $< 0.001$ & H0 отвергается: присутствует статистически значимый монотонный тренд. \\
\bottomrule
\end{tabularx}
\end{table}

Оба теста дают согласованный результат для каждой из стран. Для Германии критерий серий отвергает гипотезу о случайном чередовании наблюдений относительно медианы, а тест Kendall отвергает гипотезу об отсутствии монотонной связи со временем. Для Турции вывод полностью аналогичен: и случайность порядка наблюдений, и отсутствие монотонного тренда статистически отвергаются. Следовательно, тренд присутствует в обоих рядах. При этом сравнительный анализ графиков показывает, что в Германии он реализуется как более плавное повышение показателя, тогда как в Турции --- как более быстрый и более выраженный рост.

\subsection{Структурные сдвиги}

Визуальный анализ рядов указывает на возможное наличие структурных сдвигов, причем предполагаемые даты излома различаются по странам. На рис.~\ref{fig:germany_structural_break_section3} видно, что для Германии вероятный перелом динамики приходится на период 2008--2009 гг., после которого рост продолжается, но становится более умеренным. Для Турции на рис.~\ref{fig:turkiye_structural_break_section3} наиболее правдоподобный излом наблюдается в 2003--2004 гг., после чего повышение доли занятых в сфере услуг ускоряется.

\begin{figure}[H]
\centering
\includegraphics[width=0.88\textwidth]{figures/germany_structural_break_section3.png}
\caption{Визуальный анализ структурного сдвига для Германии}
\label{fig:germany_structural_break_section3}
\end{figure}

\begin{figure}[H]
\centering
\includegraphics[width=0.88\textwidth]{figures/turkiye_structural_break_section3.png}
\caption{Визуальный анализ структурного сдвига для Турции}
\label{fig:turkiye_structural_break_section3}
\end{figure}

Для формальной проверки используется тест Чоу. Нулевая гипотеза состоит в отсутствии структурного сдвига в линейном тренде в предполагаемой точке излома. Результаты тестирования представлены в таблице~\ref{tab:structural_breaks_section3}.

\begin{table}[H]
\centering
\begin{flushright}
\small Таблица 4
\end{flushright}
\caption{Проверка наличия структурных сдвигов показателя доли занятых в сфере услуг, \%, 1991--2025 гг.}
\label{tab:structural_breaks_section3}
\scriptsize
\begin{tabularx}{\textwidth}{@{}>{\raggedright\arraybackslash}p{0.11\textwidth} >{\centering\arraybackslash}p{0.12\textwidth} >{\raggedright\arraybackslash}p{0.22\textwidth} >{\raggedleft\arraybackslash}p{0.10\textwidth} >{\raggedleft\arraybackslash}p{0.08\textwidth} >{\raggedright\arraybackslash}X@{}}
\toprule
\makecell[tl]{\textbf{Страна}} & \makecell[c]{\textbf{Наличие} \\ \textbf{сдвига}} & \makecell[tl]{\textbf{H0}} & \makecell[tr]{\textbf{Статис-} \\ \textbf{тика}} & \makecell[tr]{\textbf{p-} \\ \textbf{value}} & \makecell[tl]{\textbf{Вывод}} \\
\midrule
Германия & Да & В линейном тренде нет структурного сдвига в 2009 году. & 200.7801 & $< 0.001$ & H0 отвергается: структурный сдвиг присутствует. \\
Турция & Да & В линейном тренде нет структурного сдвига в 2003 году. & 28.3877 & $< 0.001$ & H0 отвергается: структурный сдвиг присутствует. \\
\bottomrule
\end{tabularx}
\end{table}

Содержательная интерпретация результатов различается по странам. Для Германии структурный сдвиг естественно связывается с мировым финансовым кризисом 2008--2009 гг. и последующей перестройкой рынка труда, после которой доля занятых в услугах продолжила расти, но уже более плавно. Для Турции структурный излом около 2003 г. можно связать с трансформацией начала 2000-х годов после кризиса 2001 г., когда ускорились урбанизация и развитие торговли, туризма, финансовых и других сервисных отраслей. Тем самым в обоих случаях структурный сдвиг имеет не только статистическое, но и экономическое содержание.

Итоговый вывод по разделу состоит в том, что оба ряда имеют выраженный тренд и содержат структурные сдвиги. Германия характеризуется более плавной и зрелой траекторией роста с более поздним изломом, тогда как Турция демонстрирует более ранний структурный перелом и более интенсивное продвижение к сервисной модели занятости. Именно это различие в форме тренда и характере структурных изменений делает сопоставление двух стран особенно информативным для последующего моделирования временных рядов.

\section{Выбор и оценка ARIMA-моделей}

\subsection{Порядок интегрируемости процесса}

Для выявления стационарности рядов использованы три теста единичного корня: \textit{ADF}, \textit{KPSS} и \textit{Phillips--Perron}. Тесты \textit{ADF} и \textit{Phillips--Perron} имеют одинаковую нулевую гипотезу о наличии единичного корня, однако \textit{ADF} чувствителен к выбору лаговой структуры, тогда как \textit{Phillips--Perron} использует непараметрическую корректировку и менее зависим от спецификации ошибок. Тест \textit{KPSS}, напротив, имеет обратную нулевую гипотезу о стационарности, поэтому используется как парная проверка к двум другим тестам. Для коротких макроэкономических рядов возможны противоречия между результатами \textit{ADF}/\textit{PP} и \textit{KPSS}, особенно при наличии тренда и структурных сдвигов.

\begin{table}[H]
\centering
\begin{flushright}
\small Таблица 5
\end{flushright}
\caption{Тесты единичного корня для показателя доли занятых в сфере услуг, \%, 1991--2025 гг.}
\label{tab:integration_tests_section4}
\scriptsize
\begin{tabularx}{\textwidth}{>{\raggedright\arraybackslash}p{0.09\textwidth} >{\raggedright\arraybackslash}p{0.12\textwidth} >{\raggedright\arraybackslash}p{0.08\textwidth} >{\raggedright\arraybackslash}p{0.23\textwidth} >{\raggedleft\arraybackslash}p{0.09\textwidth} >{\raggedleft\arraybackslash}p{0.08\textwidth} >{\raggedright\arraybackslash}X}
\toprule
\makecell[tl]{\textbf{Страна}} & \makecell[tl]{\textbf{Уровень} \\ \textbf{/ разность}} & \makecell[tl]{\textbf{Тест}} & \makecell[tl]{\textbf{H0}} & \makecell[tr]{\textbf{Статис-} \\ \textbf{тика}} & \makecell[tr]{\textbf{p-} \\ \textbf{value}} & \makecell[tl]{\textbf{Вывод}} \\
\midrule
Германия & Уровень & ADF (ct) & В ряду есть единичный корень, ряд нестационарен. & -1.2264 & 0.905 & H0 не отвергается \\
Германия & Уровень & KPSS (ct) & Ряд стационарен относительно тренда. & 0.2104 & 0.012 & H0 отвергается \\
Германия & Уровень & PP (ct) & В ряду есть единичный корень, ряд нестационарен. & -1.0200 & 0.941 & H0 не отвергается \\
Германия & Первая разность & ADF (c) & В ряду есть единичный корень, ряд нестационарен. & -6.4982 & $< 0.001$ & H0 отвергается \\
Германия & Первая разность & KPSS (c) & Ряд стационарен относительно константы. & 0.6477 & 0.018 & H0 отвергается \\
Германия & Первая разность & PP (c) & В ряду есть единичный корень, ряд нестационарен. & -6.7329 & $< 0.001$ & H0 отвергается \\
Турция & Уровень & ADF (ct) & В ряду есть единичный корень, ряд нестационарен. & -1.6903 & 0.755 & H0 не отвергается \\
Турция & Уровень & KPSS (ct) & Ряд стационарен относительно тренда. & 0.1495 & 0.047 & H0 отвергается \\
Турция & Уровень & PP (ct) & В ряду есть единичный корень, ряд нестационарен. & -1.5233 & 0.821 & H0 не отвергается \\
Турция & Первая разность & ADF (c) & В ряду есть единичный корень, ряд нестационарен. & -5.4824 & $< 0.001$ & H0 отвергается \\
Турция & Первая разность & KPSS (c) & Ряд стационарен относительно константы. & 0.1364 & 0.100 & H0 не отвергается \\
Турция & Первая разность & PP (c) & В ряду есть единичный корень, ряд нестационарен. & -5.5505 & $< 0.001$ & H0 отвергается \\
\bottomrule
\end{tabularx}
\end{table}

Результаты таблицы~\ref{tab:integration_tests_section4} показывают, что для обоих рядов на уровне \textit{ADF} и \textit{Phillips--Perron} не отвергают наличие единичного корня, а \textit{KPSS} отвергает стационарность, что согласованно указывает на нестационарность уровней. После первой разности картина становится существенно ближе к стационарности. Для Турции вывод однозначен: \textit{ADF} и \textit{PP} отвергают единичный корень, а \textit{KPSS} не отвергает стационарность. Для Германии возникает частичное противоречие, поскольку \textit{ADF} и \textit{PP} указывают на стационарность первой разности, тогда как \textit{KPSS} остаётся более строгим. С учетом поведения коррелограмм и того, что вторая разность в ноутбуке выглядела избыточной, оба ряда целесообразно классифицировать как процессы $I(1)$.

Следовательно, Германия и Турция относятся к типу \textit{difference-stationary} (\textit{DS}) процессов. Стационарность достигается не простым удалением детерминированного тренда, а дифференцированием первого порядка. Поэтому в дальнейшем естественно рассматривать модели класса $ARIMA(p,1,q)$.

\subsection{Выбор параметров ARIMA}

На практике рассматривают семейства $AR$, $MA$, $ARMA$, $ARIMA$, а при наличии сезонности --- $SARIMA$. Выбор параметров обычно строится в два шага: сначала анализируются $ACF$ и $PACF$ стационарного ряда, затем оцениваются несколько близких спецификаций и сравниваются их информационные критерии и содержательная интерпретация. Поскольку оба ряда были отнесены к $I(1)$, дальнейший выбор проводится в классе моделей $ARIMA(p,1,q)$.

Для Германии анализ $ACF/PACF$ первой разности показывает лишь слабые низкоуровневые пики, поэтому нет оснований переходить к высоким порядкам модели. В качестве базовых кандидатов достаточно проверить малые значения $p$ и $q$, прежде всего $0$ или $1$. По этой причине в работе рассматриваются $ARIMA(1,1,0)$ и $ARIMA(0,1,1)$, а затем их расширения детерминированным трендом и фиктивной переменной структурного сдвига.

Для Турции первая разность выглядит еще ближе к белому шуму: и $ACF$, и $PACF$ быстро теряют выраженную структуру. Поэтому естественная отправная точка --- спецификация $ARIMA(0,1,0)$ с проверкой альтернативы $ARIMA(0,1,1)$. Такой выбор соответствует содержанию ряда: основная динамика определяется трендовой составляющей, а краткосрочная память остается короткой.

\subsection{Оценивание моделей}

Для Германии и Турции оценены по пять спецификаций: две базовые модели $ARIMA(p,1,q)$, модель с линейным трендом, модель с квадратичным трендом и модель с фиктивной переменной, соответствующей предполагаемому структурному сдвигу. Для Германии рассматривались $ARIMA(1,1,0)$, $ARIMA(0,1,1)$, а также расширения с линейным трендом, квадратичным трендом и dummy-переменной. Для Турции аналогично оценивались $ARIMA(0,1,0)$, $ARIMA(0,1,1)$ и три расширенные спецификации с линейным трендом, квадратичным трендом и dummy-переменной.

Математические формы моделей для Германии можно записать через лаговый оператор $L$ следующим образом:
\[
\begin{aligned}
ARIMA(1,1,0): \quad &(1-\phi_1L)(1-L)y_t=\varepsilon_t,\\
ARIMA(0,1,1): \quad &(1-L)y_t=(1+\theta_1L)\varepsilon_t,\\
ARIMA(1,1,0)+\text{linear trend}: \quad &(1-\phi_1L)(1-L)y_t=\beta_1 t+\varepsilon_t,\\
ARIMA(1,1,0)+\text{quadratic trend}: \quad &(1-\phi_1L)(1-L)y_t=\beta_1 t+\beta_2 t^2+\varepsilon_t,\\
ARIMA(1,1,0)+\text{dummy}: \quad &(1-\phi_1L)(1-L)y_t=\delta D_t+\varepsilon_t.
\end{aligned}
\]

Для Турции соответствующие модели имеют вид
\[
\begin{aligned}
ARIMA(0,1,0): \quad &(1-L)y_t=\varepsilon_t,\\
ARIMA(0,1,1): \quad &(1-L)y_t=(1+\theta_1L)\varepsilon_t,\\
ARIMA(0,1,0)+\text{linear trend}: \quad &(1-L)y_t=\beta_1 t+\varepsilon_t,\\
ARIMA(0,1,0)+\text{quadratic trend}: \quad &(1-L)y_t=\beta_1 t+\beta_2 t^2+\varepsilon_t,\\
ARIMA(0,1,0)+\text{dummy}: \quad &(1-L)y_t=\delta D_t+\varepsilon_t.
\end{aligned}
\]

Оценки коэффициентов и их значимость для Германии представлены в таблице~\ref{tab:germany_coeff_section4}. Наиболее содержательно сильной выглядит модель с квадратичным трендом: в ней значимы линейный и квадратичный тренд, а также авторегрессионный коэффициент. В модели с линейным трендом значим тренд, но коэффициент при $AR(1)$ незначим. В dummy-спецификации фиктивная переменная незначима.

\begin{table}[H]
\centering
\begin{flushright}
\small Таблица 6
\end{flushright}
\caption{Оценки коэффициентов моделей ряда доли занятых в сфере услуг для Германии, \%, 1991--2025 гг.}
\label{tab:germany_coeff_section4}
\scriptsize
\begin{tabularx}{\textwidth}{>{\raggedright\arraybackslash}p{0.24\textwidth} >{\raggedright\arraybackslash}p{0.17\textwidth} >{\raggedleft\arraybackslash}p{0.11\textwidth} >{\raggedleft\arraybackslash}p{0.11\textwidth} >{\raggedright\arraybackslash}X}
\toprule
\textbf{Модель} & \textbf{Параметр} & \textbf{Оценка} & \textbf{p-value} & \textbf{Значимость} \\
\midrule
\makecell[tl]{ARIMA(1,1,0)} & ar.L1 & 0.3677 & 0.067 & не значим \\
\makecell[tl]{ARIMA(1,1,0)} & sigma2 & 0.3718 & $< 0.001$ & значим \\
\makecell[tl]{ARIMA(0,1,1)} & ma.L1 & 0.2295 & 0.332 & не значим \\
\makecell[tl]{ARIMA(0,1,1)} & sigma2 & 0.3943 & $< 0.001$ & значим \\
\makecell[tl]{ARIMA(1,1,0) \\ + linear trend} & trend\_linear & 0.4411 & $< 0.001$ & значим \\
\makecell[tl]{ARIMA(1,1,0) \\ + linear trend} & ar.L1 & -0.1538 & 0.419 & не значим \\
\makecell[tl]{ARIMA(1,1,0) \\ + linear trend} & sigma2 & 0.2309 & $< 0.001$ & значим \\
\makecell[tl]{ARIMA(1,1,0) \\ + quadratic trend} & trend\_linear & 0.8343 & $< 0.001$ & значим \\
\makecell[tl]{ARIMA(1,1,0) \\ + quadratic trend} & trend\_quadratic & -0.0108 & 0.007 & значим \\
\makecell[tl]{ARIMA(1,1,0) \\ + quadratic trend} & ar.L1 & -0.3991 & 0.002 & значим \\
\makecell[tl]{ARIMA(1,1,0) \\ + quadratic trend} & sigma2 & 0.1631 & $< 0.001$ & значим \\
\makecell[tl]{ARIMA(1,1,0) \\ + dummy} & shift\_dummy & -0.1050 & 0.984 & не значим \\
\makecell[tl]{ARIMA(1,1,0) \\ + dummy} & ar.L1 & 0.3778 & 0.061 & не значим \\
\makecell[tl]{ARIMA(1,1,0) \\ + dummy} & sigma2 & 0.3714 & $< 0.001$ & значим \\
\bottomrule
\end{tabularx}
\end{table}

Оценки для Турции приведены в таблице~\ref{tab:turkiye_coeff_section4}. Здесь особенно сильной выглядит модель с линейным трендом: коэффициент тренда значим, а спецификация в целом дает наилучшие информационные критерии. В модели с квадратичным трендом квадратичный член оказывается незначимым, а в модели с dummy фиктивная переменная незначима.

\begin{table}[H]
\centering
\begin{flushright}
\small Таблица 7
\end{flushright}
\caption{Оценки коэффициентов моделей ряда доли занятых в сфере услуг для Турции, \%, 1991--2025 гг.}
\label{tab:turkiye_coeff_section4}
\scriptsize
\begin{tabularx}{\textwidth}{>{\raggedright\arraybackslash}p{0.24\textwidth} >{\raggedright\arraybackslash}p{0.17\textwidth} >{\raggedleft\arraybackslash}p{0.11\textwidth} >{\raggedleft\arraybackslash}p{0.11\textwidth} >{\raggedright\arraybackslash}X}
\toprule
\textbf{Модель} & \textbf{Параметр} & \textbf{Оценка} & \textbf{p-value} & \textbf{Значимость} \\
\midrule
\makecell[tl]{ARIMA(0,1,0)} & sigma2 & 1.5102 & 0.001 & значим \\
\makecell[tl]{ARIMA(0,1,1)} & ma.L1 & 0.2686 & 0.141 & не значим \\
\makecell[tl]{ARIMA(0,1,1)} & sigma2 & 1.2693 & $< 0.001$ & значим \\
\makecell[tl]{ARIMA(0,1,0) \\ + linear trend} & trend\_linear & 0.7867 & $< 0.001$ & значим \\
\makecell[tl]{ARIMA(0,1,0) \\ + linear trend} & sigma2 & 0.8913 & $< 0.001$ & значим \\
\makecell[tl]{ARIMA(0,1,0) \\ + quadratic trend} & trend\_linear & 0.9892 & 0.009 & значим \\
\makecell[tl]{ARIMA(0,1,0) \\ + quadratic trend} & trend\_quadratic & -0.0055 & 0.613 & не значим \\
\makecell[tl]{ARIMA(0,1,0) \\ + quadratic trend} & sigma2 & 0.8805 & $< 0.001$ & значим \\
\makecell[tl]{ARIMA(0,1,0) \\ + dummy} & shift\_dummy & 1.2925 & 1.000 & не значим \\
\makecell[tl]{ARIMA(0,1,0) \\ + dummy} & sigma2 & 1.4596 & $< 0.001$ & значим \\
\bottomrule
\end{tabularx}
\end{table}

Сравнение моделей по информационным критериям, а также по выполнению условий стационарности и обратимости, приведено в таблице~\ref{tab:model_compare_section4}. Для Германии наилучшие значения $AIC$ и $BIC$ имеет модель $ARIMA(1,1,0)+$quadratic trend. Для Турции наилучшей по тем же критериям оказывается $ARIMA(0,1,0)+$linear trend. Тем самым обе страны требуют учета детерминированной трендовой составляющей, однако для Германии оправдан более гибкий квадратичный тренд, тогда как для Турции достаточно линейного.

\begin{table}[H]
\centering
\begin{flushright}
\small Таблица 8
\end{flushright}
\caption{Сравнение оцененных моделей ряда доли занятых в сфере услуг, \%, 1991--2025 гг.}
\label{tab:model_compare_section4}
\scriptsize
\begin{tabularx}{\textwidth}{>{\raggedright\arraybackslash}p{0.08\textwidth} >{\raggedright\arraybackslash}p{0.17\textwidth} >{\centering\arraybackslash}p{0.09\textwidth} >{\raggedleft\arraybackslash}p{0.07\textwidth} >{\raggedleft\arraybackslash}p{0.07\textwidth} >{\raggedright\arraybackslash}p{0.16\textwidth} >{\raggedright\arraybackslash}X}
\toprule
\textbf{Страна} & \textbf{Модель} & \textbf{(p,d,q)} & \textbf{AIC} & \textbf{BIC} & \makecell[tl]{\textbf{Стационарность} \\ \textbf{/ обратимость}} & \textbf{Вывод} \\
\midrule
Германия & \makecell[tl]{ARIMA(1,1,0)} & (1,1,0) & 65.00 & 67.99 & стационарность: выполнена; обратимость: неприменимо & резервная модель \\
Германия & \makecell[tl]{ARIMA(0,1,1)} & (0,1,1) & 65.04 & 67.97 & стационарность: неприменимо; обратимость: выполнена & резервная модель \\
Германия & \makecell[tl]{ARIMA(1,1,0) \\ + linear trend} & (1,1,0) & 51.28 & 55.77 & стационарность: выполнена; обратимость: неприменимо & содержательно полезная спецификация \\
Германия & \makecell[tl]{ARIMA(1,1,0) \\ + quadratic trend} & (1,1,0) & 41.80 & 47.79 & стационарность: выполнена; обратимость: неприменимо & одна из сильнейших по $AIC/BIC$ \\
Германия & \makecell[tl]{ARIMA(1,1,0) \\ + dummy} & (1,1,0) & 66.97 & 71.46 & стационарность: выполнена; обратимость: неприменимо & резервная модель \\
Турция & \makecell[tl]{ARIMA(0,1,0)} & (0,1,0) & 109.25 & 110.75 & стационарность: неприменимо; обратимость: неприменимо & резервная модель \\
Турция & \makecell[tl]{ARIMA(0,1,1)} & (0,1,1) & 102.45 & 105.38 & стационарность: неприменимо; обратимость: выполнена & резервная модель \\
Турция & \makecell[tl]{ARIMA(0,1,0) \\ + linear trend} & (0,1,0) & 93.85 & 96.85 & стационарность: неприменимо; обратимость: неприменимо & одна из сильнейших по $AIC/BIC$ \\
Турция & \makecell[tl]{ARIMA(0,1,0) \\ + quadratic trend} & (0,1,0) & 95.45 & 99.94 & стационарность: неприменимо; обратимость: неприменимо & содержательно полезная спецификация \\
Турция & \makecell[tl]{ARIMA(0,1,0) \\ + dummy} & (0,1,0) & 110.13 & 113.12 & стационарность: неприменимо; обратимость: неприменимо & резервная модель \\
\bottomrule
\end{tabularx}
\end{table}

\subsection{Анализ остатков}

Для диагностики остатков использованы тесты \textit{Ljung--Box} на автокорреляцию и \textit{Jarque--Bera} на нормальность. Ниже результаты представлены в формате итогового сравнения моделей для каждой страны, где одновременно учитываются стационарность и обратимость, значения $AIC/BIC$, свойства остатков и итоговый выбор лучшей спецификации.

\begin{table}[H]
\centering
\begin{flushright}
\small Таблица 9
\end{flushright}
\caption{Сравнение моделей и диагностика остатков для ряда доли занятых в сфере услуг в Германии, \%, 1991--2025 гг.}
\label{tab:residuals_germany_section4}
\scriptsize
{\setlength{\tabcolsep}{3pt}
\resizebox{\textwidth}{!}{%
\begin{tabularx}{\textwidth}{@{}>{\raggedright\arraybackslash}p{0.10\textwidth} >{\raggedright\arraybackslash}X >{\raggedright\arraybackslash}p{0.12\textwidth} >{\raggedright\arraybackslash}X >{\raggedright\arraybackslash}p{0.15\textwidth}@{}}
\toprule
\makecell[c]{\textbf{Модель} \\ \textbf{ARIMA}} &
\makecell[c]{\textbf{Стационарность} \\ \textbf{/ обратимость}} &
\makecell[c]{\textbf{Ошибка модели,} \\ \textbf{инф.\ критерии}} &
\makecell[c]{\textbf{Анализ остатков} \\ \textbf{(автокорреляция,} \\ \textbf{нормальность)}} &
\makecell[c]{\textbf{Наилучшая} \\ \textbf{модель}} \\
\midrule
\makecell[tl]{ARIMA(1,1,0)} &
\makecell[tl]{стац.: выполнена; \\ обр.: неприменимо} &
\makecell[tl]{AIC = 65.00; \\ BIC = 67.99} &
\makecell[tl]{автокорр.: нет; \\ норм.: не вып.} &
нет \\
\makecell[tl]{ARIMA(0,1,1)} &
\makecell[tl]{стац.: неприменимо; \\ обр.: выполнена} &
\makecell[tl]{AIC = 65.04; \\ BIC = 67.97} &
\makecell[tl]{автокорр.: нет; \\ норм.: не вып.} &
нет \\
\makecell[tl]{ARIMA(1,1,0) \\ + linear trend} &
\makecell[tl]{стац.: выполнена; \\ обр.: неприменимо} &
\makecell[tl]{AIC = 51.28; \\ BIC = 55.77} &
\makecell[tl]{автокорр.: нет; \\ норм.: не вып.} &
нет \\
\makecell[tl]{ARIMA(1,1,0) \\ + quadratic trend} &
\makecell[tl]{стац.: выполнена; \\ обр.: неприменимо} &
\makecell[tl]{AIC = 41.80; \\ BIC = 47.79} &
\makecell[tl]{автокорр.: нет; \\ норм.: не вып.} &
да \\
\makecell[tl]{ARIMA(1,1,0) \\ + dummy} &
\makecell[tl]{стац.: выполнена; \\ обр.: неприменимо} &
\makecell[tl]{AIC = 66.97; \\ BIC = 71.46} &
\makecell[tl]{автокорр.: нет; \\ норм.: не вып.} &
нет \\
\bottomrule
\end{tabularx}
}
}
\end{table}

\begin{table}[H]
\centering
\begin{flushright}
\small Таблица 10
\end{flushright}
\caption{Сравнение моделей и диагностика остатков для ряда доли занятых в сфере услуг в Турции, \%, 1991--2025 гг.}
\label{tab:residuals_turkiye_section4}
\scriptsize
{\setlength{\tabcolsep}{3pt}
\resizebox{\textwidth}{!}{%
\begin{tabularx}{\textwidth}{@{}>{\raggedright\arraybackslash}p{0.10\textwidth} >{\raggedright\arraybackslash}X >{\raggedright\arraybackslash}p{0.12\textwidth} >{\raggedright\arraybackslash}X >{\raggedright\arraybackslash}p{0.15\textwidth}@{}}
\toprule
\makecell[c]{\textbf{Модель} \\ \textbf{ARIMA}} &
\makecell[c]{\textbf{Стационарность} \\ \textbf{/ обратимость}} &
\makecell[c]{\textbf{Ошибка модели,} \\ \textbf{инф.\ критерии}} &
\makecell[c]{\textbf{Анализ остатков} \\ \textbf{(автокорреляция,} \\ \textbf{нормальность)}} &
\makecell[c]{\textbf{Наилучшая} \\ \textbf{модель}} \\
\midrule
\makecell[tl]{ARIMA(0,1,0)} &
\makecell[tl]{стац.: неприменимо; \\ обр.: неприменимо} &
\makecell[tl]{AIC = 109.25; \\ BIC = 110.75} &
\makecell[tl]{автокорр.: нет; \\ норм.: не вып.} &
нет \\
\makecell[tl]{ARIMA(0,1,1)} &
\makecell[tl]{стац.: неприменимо; \\ обр.: выполнена} &
\makecell[tl]{AIC = 102.45; \\ BIC = 105.38} &
\makecell[tl]{автокорр.: нет; \\ норм.: не вып.} &
нет \\
\makecell[tl]{ARIMA(0,1,0) \\ + linear trend} &
\makecell[tl]{стац.: неприменимо; \\ обр.: неприменимо} &
\makecell[tl]{AIC = 93.85; \\ BIC = 96.85} &
\makecell[tl]{автокорр.: нет; \\ норм.: не вып.} &
да \\
\makecell[tl]{ARIMA(0,1,0) \\ + quadratic trend} &
\makecell[tl]{стац.: неприменимо; \\ обр.: неприменимо} &
\makecell[tl]{AIC = 95.45; \\ BIC = 99.94} &
\makecell[tl]{автокорр.: нет; \\ норм.: не вып.} &
нет \\
\makecell[tl]{ARIMA(0,1,0) \\ + dummy} &
\makecell[tl]{стац.: неприменимо; \\ обр.: неприменимо} &
\makecell[tl]{AIC = 110.13; \\ BIC = 113.12} &
\makecell[tl]{автокорр.: нет; \\ норм.: не вып.} &
нет \\
\bottomrule
\end{tabularx}
}
}
\end{table}

Для Германии в качестве наилучшей выбрана модель $ARIMA(1,1,0)$ с квадратичным трендом, поскольку именно она даёт минимальные значения $AIC$ и $BIC$ при отсутствии автокорреляции остатков. Для Турции лучшей является модель $ARIMA(0,1,0)$ с линейным трендом, которая также сочетает наилучшие информационные критерии с приемлемой диагностикой остатков. В обеих странах нормальность остатков не выполняется, поэтому выбор модели основан прежде всего на сочетании $AIC/BIC$ и отсутствия остаточной автокорреляции. Тем самым спецификация для Германии оказывается более гибкой по тренду, тогда как для Турции достаточной является более простая линейная форма.

\subsection{Прогноз}

Для выбранных моделей построен прогноз на шесть шагов вперед; соответствующие точечные значения и 95\%-ные доверительные интервалы представлены в таблице~\ref{tab:forecast_section4}. Графическое представление исходных рядов, подогнанных значений и прогноза приведено на рис.~\ref{fig:forecast_germany_section4} и~\ref{fig:forecast_turkiye_section4}.

\begin{table}[H]
\centering
\begin{flushright}
\small Таблица 11
\end{flushright}
\caption{Прогноз показателя доли занятых в сфере услуг на 2026--2031 гг., \%.}
\label{tab:forecast_section4}
\scriptsize
\begin{tabularx}{\textwidth}{>{\raggedright\arraybackslash}p{0.10\textwidth} >{\raggedleft\arraybackslash}p{0.07\textwidth} >{\raggedleft\arraybackslash}p{0.07\textwidth} >{\raggedleft\arraybackslash}p{0.16\textwidth} >{\raggedleft\arraybackslash}p{0.18\textwidth} >{\raggedleft\arraybackslash}X}
\toprule
\textbf{Страна} & \textbf{Год} & \textbf{Шаг} & \makecell[tr]{\textbf{Точечный} \\ \textbf{прогноз}} & \makecell[tr]{\textbf{Нижняя граница} \\ \textbf{95\%}} & \makecell[tr]{\textbf{Верхняя граница} \\ \textbf{95\%}} \\
\midrule
Германия & 2026 & 1 & 72.9254 & 72.1339 & 73.7168 \\
Германия & 2027 & 2 & 73.0126 & 72.0893 & 73.9359 \\
Германия & 2028 & 3 & 73.0231 & 71.9211 & 74.1251 \\
Германия & 2029 & 4 & 73.0340 & 71.8018 & 74.2663 \\
Германия & 2030 & 5 & 73.0147 & 71.6564 & 74.3729 \\
Германия & 2031 & 6 & 72.9773 & 71.5068 & 74.4477 \\
Турция & 2026 & 1 & 59.2762 & 57.4258 & 61.1265 \\
Турция & 2027 & 2 & 60.0628 & 57.4460 & 62.6797 \\
Турция & 2028 & 3 & 60.8495 & 57.6446 & 64.0545 \\
Турция & 2029 & 4 & 61.6362 & 57.9355 & 65.3370 \\
Турция & 2030 & 5 & 62.4229 & 58.2853 & 66.5605 \\
Турция & 2031 & 6 & 63.2096 & 58.6771 & 67.7421 \\
\bottomrule
\end{tabularx}
\end{table}

\begin{figure}[H]
\centering
\includegraphics[width=0.88\textwidth]{figures/forecast_germany_section4.png}
\caption{Прогноз для Германии: исходный ряд, fitted values, forecast и 95\% confidence interval}
\label{fig:forecast_germany_section4}
\end{figure}

\begin{figure}[H]
\centering
\includegraphics[width=0.88\textwidth]{figures/forecast_turkiye_section4.png}
\caption{Прогноз для Турции: исходный ряд, fitted values, forecast и 95\% confidence interval}
\label{fig:forecast_turkiye_section4}
\end{figure}

Прогнозные траектории двух стран различаются достаточно четко. Для Германии ожидается сохранение высокого уровня показателя при очень плавной динамике: прогноз практически стабилизируется около 73\% общей занятости. Для Турции прогноз остается восходящим и предполагает дальнейший рост от 59.28\% в 2026 г. до 63.21\% в 2031 г. Следовательно, в краткосрочном горизонте Турция продолжает расти быстрее, тогда как Германия выглядит более стабильной и инерционной.

\subsection{Качество прогноза}

\begin{table}[H]
\centering
\begin{flushright}
\small Таблица 12
\end{flushright}
\caption{Метрики качества прогноза для показателя доли занятых в сфере услуг, \%, 1991--2025 гг.}
\label{tab:forecast_quality_section4}
\begin{tabular}{lrrr}
\toprule
\textbf{Страна} & \textbf{MAE} & \textbf{RMSE} & \textbf{MAPE} \\
\midrule
Германия & 2.2472 & 9.8399 & 3.7961 \\
Турция & 1.6175 & 5.2111 & 4.4778 \\
\bottomrule
\end{tabular}
\end{table}

Согласно таблице~\ref{tab:forecast_quality_section4}, по относительной ошибке прогноз точнее для Германии: значение $MAPE$ составляет 3.7961 против 4.4778 для Турции. Абсолютные метрики $MAE$ и $RMSE$ ниже у Турции, однако их интерпретацию следует вести с учетом масштаба ряда. В совокупности результаты указывают на приемлемое краткосрочное качество обеих моделей, но надежность прогноза следует оценивать как умеренную: отсутствие автокорреляции остатков поддерживает адекватность моделей, тогда как ненормальность остатков ограничивает строгость интерпретации доверительных интервалов.

\section{Заключение}

В работе исследованы временные ряды показателя \textit{Employment in services (\% of total employment)} для Германии и Турции за 1991--2025 гг. На уровне описательной статистики и графического анализа установлено, что Германия на всем интервале характеризуется более высоким уровнем занятости в сфере услуг, тогда как Турция демонстрирует более быстрый рост показателя. В обоих рядах выявлены выраженные трендовые компоненты, статистически подтвержденные тестами на отсутствие тренда, а также структурные сдвиги: для Германии около 2009 г., для Турции около 2003 г. Проверка стационарности показала, что оба процесса следует трактовать как $I(1)$ и \textit{DS}-процессы, что сделало естественным переход к моделям класса $ARIMA(p,1,q)$.

По итогам сравнительного оценивания и диагностики остатков в качестве рабочих спецификаций выбраны модель с квадратичным трендом для Германии и модель с линейным трендом для Турции. Прогнозирование на шесть шагов вперед показало, что для Германии наиболее вероятна стабилизация показателя около достигнутого высокого уровня, тогда как для Турции сохраняется более выраженная восходящая динамика. В целом сопоставление двух стран позволяет увидеть различие между более зрелой и плавной траекторией сервисизации занятости в Германии и более интенсивной догоняющей трансформацией в Турции. Именно это различие и определяет содержательную ценность проведенного анализа.

\section{Список литературы}

\begin{enumerate}
\item Shumway R.H., Stoffer D.S. \textit{Time Series Analysis and Its Applications: With R Examples}. Fourth Edition, 2017.
\item Лекции курса «Временные ряды и их практическое применение».
\end{enumerate}

\section{Приложение}

В приложении приведены дополнительные результаты, полученные в Python и не включённые в основной текст отчёта: детализированные проверки стационарности и сводная диагностика остатков моделей.

\begin{figure}[H]
\centering
\includegraphics[width=0.92\textwidth]{figures/appendix_adf_results.png}
\caption{Результаты теста ADF для рядов Германии и Турции на уровне и в первой разности}
\label{fig:appendix_adf}
\end{figure}

\begin{figure}[H]
\centering
\includegraphics[width=0.92\textwidth]{figures/appendix_kpss_results.png}
\caption{Результаты теста KPSS для рядов Германии и Турции на уровне и в первой разности}
\label{fig:appendix_kpss}
\end{figure}

\begin{figure}[H]
\centering
\includegraphics[width=0.92\textwidth]{figures/appendix_pp_results.png}
\caption{Результаты теста Phillips--Perron для рядов Германии и Турции на уровне и в первой разности}
\label{fig:appendix_pp}
\end{figure}

\begin{figure}[H]
\centering
\includegraphics[width=0.92\textwidth]{figures/appendix_residual_diagnostics.png}
\caption{Результаты тестов Ljung--Box и Jarque--Bera в диагностике остатков моделей}
\label{fig:appendix_residuals}
\end{figure}

\end{document}
